% !PRINTED 224
\weberpage{224}
%<<CONT
$n$ aufgehende Primzahl, so nehme man zwei Wurzeln $x_{\mu,\mu'}$, $x_{\nu,\nu'}$
in dieselbe oder in verschiedene Reihen auf, je nachdem

\[
\mu\nu' - \nu\mu' \equiv 0 \ (\mathrm{mod}\ p)
\]

oder nicht; gehören hiernach $x_{\nu,\nu'}$ und $x_{\nu_1,\nu_1'}$ in eine Reihe mit
$x_{\mu,\mu'}$, so gehören sie auch untereinander in dieselbe Reihe. Die
Anzahl der so gebildeten Reihen ist, wie leicht nachzuweisen,
$p + 1$ und man erhält auf diese Weise als erste Resolvente der
Teilungsgleichung eine zum Divisor $p$ gehörige Transformationsgleichung.
Wir werden später auf anderem Wege zeigen, wie die
Transformationsgleichungen für zusammengesetzte Divisoren auf
solche für Primzahldivisoren zurückgeführt werden können.

\begin{center}
\rule{3cm}{0.4pt}
\end{center}
